\documentclass{article}
\usepackage[T1]{fontenc}
\usepackage{lmodern}
\usepackage{graphicx}
\usepackage{amsmath,amssymb}
\usepackage[a4paper,margin=2cm]{geometry}
\usepackage[absolute,overlay]{textpos}
\setlength{\TPHorizModule}{1mm}
\setlength{\TPVertModule}{1mm}
\usepackage[indonesian]{babel}
\usepackage{hyperref}
\setlength{\emergencystretch}{2em}
% unit: Halaman 22

\begin{document}

% === Halaman 22 (python/native) ===

Rinaldi Munir/II4021 Kriptografi

x13 + x11 + x9 + x8 + x6 + x5 + x4 + x3 +  1 

x8 +  x4 +  x3 + x + 1 

x5 + x3 + 


x13               + x9 + x8 + x6 + x5 

+

x11 +                            + x4 + x3 +  1 

x11 +                x7 + x6 + x4 + x3  

+

x7 + x6                 +   1 

Sisa pembagian

Jadi, x13 + x11 + x9 + x8 + x6 + x5 + x4 + x3 + 1 (mod f(x)) = x7 + x6 + 1  = 11000001

22

\end{document}
